\documentclass{article}

% V0.1 uses a preprint layout. Do not change this to a conference track until
% the submission venue and author list have been confirmed.
\usepackage[preprint]{neurips_2026}
\usepackage[utf8]{inputenc}
\usepackage[T1]{fontenc}
\usepackage{courier}
\usepackage{hyperref}
\usepackage{url}
\usepackage{booktabs}
\usepackage{amsmath,amsfonts,amssymb}
\usepackage{microtype}
\usepackage{xcolor}
\usepackage{enumitem}

\newcommand{\benchmark}{\textsc{EviSuff-Finance}}
\newcommand{\mses}{\textsc{MSES}}
\newcommand{\pbs}{\textsc{PBS}}
\newcommand{\upr}{\textsc{UPR}}
\newcommand{\cas}{\textsc{CAS}}
\hypersetup{colorlinks=true,linkcolor=black,citecolor=black,urlcolor=blue}

\title{When Relevant Evidence Is Not Enough: Counterfactual Evidence-Sufficiency Evaluation for Long-Document Financial Agents}

% Replace only after authorship and affiliations are jointly confirmed.
\author{Anonymous Author(s)\\
Affiliation to be confirmed\\
\texttt{contact to be confirmed}}

\begin{document}
\maketitle

\begin{abstract}
Long-document financial agents are increasingly evaluated by answer accuracy,
retrieval scores, or document-level citations. These signals do not establish
whether an agent changes its behavior when the evidence required for an answer
is absent. We introduce \benchmark{}, a counterfactual evaluation protocol for
financial filings. Each question is annotated with one or more minimal
sufficient evidence sets (\mses{}). A paired intervention removes a smallest
hitting set that intersects every \mses{}, making the question unanswerable
from the allowed corpus while preserving most topical context. Agents should
answer under sufficient evidence, abstain after necessary-evidence removal, and
remain stable under an irrelevant-removal control. This V0.1 paper contributes
(i) the task definition and annotation protocol, (ii) an open implementation
of the intervention and paired metrics, and (iii) a reproducible audit of the
public IPO Finance Agent repository. The audit confirms a useful 70-question
IPO seed set but also shows why claim-to-passage human provenance is necessary
for the proposed evaluation. We report no model-comparison results: the
human-validated benchmark and preregistered experiment remain future work.
\end{abstract}

\section{Introduction}

Financial agents increasingly operate over lengthy regulatory filings,
combining tables, footnotes, risk disclosures, accounting recasts, and
governance information. Correctness alone is not enough in this setting: a
response can look plausible, cite a whole filing, and still be unsupported by
the evidence made available to the model. The distinction matters for research
agents, due diligence assistants, and systems used to support high-stakes
professional workflows.

Recent benchmarks document important limitations of financial retrieval and
reasoning. IPO Finance Agent extends financial-agent evaluation to IPO
diligence and long S-1 filings \citep{benhenda2026ipo}. Fin-RATE separates
ground-truth-context and retrieval-augmented settings over SEC tasks
\citep{jiang2026finrate}; LOFin/HiREC studies evidence curation over large SEC
corpora \citep{choe2025hirec}; and FinAgentBench assesses agentic document and
passage selection \citep{choi2025finagentbench}. In general RAG, sufficient
context \citep{joren2024sufficient}, answerable/unanswerable search
\citep{xie2026oversearch}, citation evaluation \citep{gao2023alce}, and
conflicting evidence \citep{lee2025magic,niu2026arbgraph} have each received
direct attention. Deep-research benchmarks further expose difficulties in
evidence-backed multi-step synthesis \citep{sharma2025researchrubrics,gou2025mind2web}.

These works leave a precise diagnostic question open: \emph{does a financial
agent causally track the evidence necessary for its answer?} An answerable
context can contain relevant passages without being sufficient, while removing
one passage from a multi-source answer may leave an alternative complete proof
path intact. We therefore require annotators to enumerate multiple minimal
sufficient evidence sets and construct a removal that breaks all known paths.

This paper is deliberately modest. It is a V0.1 protocol and reproducibility
audit, not a completed empirical benchmark. Its contributions are:

\begin{enumerate}[leftmargin=1.3em]
    \item a formal paired intervention for evidence sufficiency based on
    human-verified \mses{} and minimum hitting sets;
    \item an annotation, evaluation, and statistical protocol for long-document
    financial agents;
    \item a tested open-source implementation of the core intervention and
    metrics; and
    \item a scoped, reproducible audit of public IPO Finance Agent artifacts to
    motivate the need for claim-level evidence provenance.
\end{enumerate}

\section{Related work}

\paragraph{Financial documents and agents.}
IPO Finance Agent introduces a public set of 70 SpaceX IPO questions and a
contextual-retrieval harness designed for long S-1 documents
\citep{benhenda2026ipo}. Fin-RATE evaluates financial analytics across single
documents, longitudinal data, and cross-entity comparisons
\citep{jiang2026finrate}. LOFin/HiREC focuses on hierarchical retrieval and
evidence curation for standardized SEC documents \citep{choe2025hirec}, while
FinSage and FinAgentBench target financial filing QA and agentic retrieval
\citep{wang2025finsage,choi2025finagentbench}. \benchmark{} uses such tasks as
a domain but does not propose a competing retrieval method.

\paragraph{Sufficiency, abstention, and citations.}
Sufficient Context distinguishes errors caused by insufficient retrieval from
those caused by failures to use adequate evidence \citep{joren2024sufficient}.
Over-Searching shows that noisy search can harm abstention on unanswerable
queries \citep{xie2026oversearch}, and ERA studies evidence-based reliability
alignment for abstention \citep{shin2026era}. ALCE formalizes citation
correctness and completeness \citep{gao2023alce}; more recent work evaluates
source attribution in deep-research reports \citep{onweller2026cited}. Our
focus differs: we assess whether answer and abstention behavior changes after
a controlled removal of necessary evidence.

\paragraph{Conflicts, time, and research agents.}
MAGIC and ArbGraph address contradictory or conflicting retrieved evidence
\citep{lee2025magic,niu2026arbgraph}. Look-Ahead-Bench studies temporal
look-ahead bias in financial LLMs \citep{benhenda2026lookahead}. ResearchRubrics
and Mind2Web~2 evaluate long-horizon research agents and their evidence-backed
outputs \citep{sharma2025researchrubrics,gou2025mind2web}. Contradiction and
temporal interventions are important extensions, but \benchmark{} first isolates
the simpler causal question of whether evidence is sufficient at all.

\section{Task definition}

For a question $q$, let $D$ be a fixed corpus of filings and let $E$ be its
set of stable evidence units. A unit is a retrievable passage or table region
with a document ID, section, text, source URL, and optional publication date.
Let $\mathcal{M}_q = \{M_1, \ldots, M_k\}$ be human-validated minimal
sufficient evidence sets. An available context $C \subseteq E$ is answerable
when at least one complete set remains:

\begin{equation}
    \operatorname{answerable}(q,C)=1\quad \Longleftrightarrow\quad
    \exists M_i \in \mathcal{M}_q \text{ such that } M_i \subseteq C.
\end{equation}

The necessary-removal intervention selects a minimum hitting set:

\begin{equation}
H^* \in \arg\min_{H \subseteq \cup_i M_i} |H|
\quad \text{s.t.} \quad H \cap M_i \neq \varnothing\;\; \forall i.
\end{equation}

The context $C^- = C \setminus H^*$ contains no known complete \mses{} and is
labeled unanswerable with respect to the allowed corpus. This definition is
stronger than deleting a random supporting passage: an alternative evidence
path cannot silently survive the intervention.

Each item produces four paired conditions:

\begin{enumerate}[leftmargin=1.3em]
    \item \texttt{full}: retrieved or curated context with sufficient evidence;
    \item \texttt{gold\_only}: a single complete \mses{} without distractors;
    \item \texttt{necessary\_removal}: the hitting set is removed; and
    \item \texttt{irrelevant\_removal}: an irrelevant distractor is removed.
\end{enumerate}

The prompt requires a structured response with an answerability decision,
probability of answerability, answer, atomic claims, passage IDs cited for each
claim, and a missing-information explanation when abstaining. The model is not
told which condition generated the context.

\section{Benchmark construction protocol}

The planned full benchmark contains 60 questions (minimum 40) drawn from at
least three, preferably five, S-1/F-1 filings. Every source is frozen by SEC
accession number, filing date, URL, hash, and extraction version. The initial
question pool can include public IPO Finance Agent items, but the test set must
not be dominated by one issuer: SpaceX should contribute fewer than 40\% of
the final items.

Each record specifies a reference answer decomposed into atomic claims,
evidence units with stable passage identifiers, one or more \mses{},
distractors, an answerability label, cutoff-date metadata, and annotator and
adjudication information. At least 25\% of the full set is double annotated;
the 20--30-item minimum viable benchmark is double annotated in full. We will
report Cohen's $\kappa$ for answerability and evidence-unit precision, recall,
and F1 for spans and sets. A two-person review verifies every
\texttt{necessary\_removal} context.

Disclosure-of-absence questions require special care. ``The filing does not
disclose X'' cannot be supported by one missing string. Annotators must fix a
bounded search scope, list relevant sections and synonyms, and document the
verification procedure. These items will be analyzed separately from
affirmative-evidence questions.

\section{Evaluation}

\paragraph{Primary paired metrics.}
For $N$ question-level pairs, paired behavior success measures whether a model
answers correctly in the full context and abstains after necessary removal:

\begin{equation}
\pbs=\frac{1}{N}\sum_{j=1}^{N}\mathbb{1}[\mathrm{correct}_{j,\mathrm{full}}
\wedge \neg\mathrm{abstain}_{j,\mathrm{full}}
\wedge \mathrm{abstain}_{j,\mathrm{remove}}].
\end{equation}

Unsupported persistence rate is the fraction of necessary-removal cases in
which the model still answers:

\begin{equation}
\upr=\frac{1}{N}\sum_{j=1}^{N}\mathbb{1}[\neg\mathrm{abstain}_{j,\mathrm{remove}}].
\end{equation}

Counterfactual abstention shift is
$\cas=P(\mathrm{abstain}\mid\mathrm{necessary\ removal})-
P(\mathrm{abstain}\mid\mathrm{full})$. A system that tracks necessity should
have high \pbs{}, low \upr{}, positive \cas{}, and high stability under
\texttt{irrelevant\_removal}.

\paragraph{Secondary metrics.}
Retrieval is evaluated with Recall@k, nDCG@k, and \emph{Complete Evidence
Recall}: whether at least one entire \mses{} was retrieved. This avoids
over-crediting a retriever that retrieves individual relevant passages but
misses a necessary multi-passage chain. Generation evaluation includes atomic
claim precision/recall, citation precision, citation completeness, selective
accuracy, risk--coverage curves, Brier score for answerability probability,
latency, tokens, tool calls, and cost.

\paragraph{Statistical analysis.}
We will bootstrap at the question level to preserve paired conditions, report
95\% confidence intervals, use McNemar tests for matched behavior changes, and
fit mixed-effects logistic models with model and question/filing effects where
appropriate. Holm correction applies to secondary families of hypotheses.
LLM judging is auxiliary until calibrated against at least 150
human-labeled claim--citation pairs, with judge--human agreement and an error
analysis reported.

\section{Implementation and planned baselines}

The accompanying package validates JSONL annotations, computes a minimum
hitting set, materializes all four conditions, calculates paired metrics, and
provides a question-level bootstrap. Unit tests cover multiple evidence paths,
condition construction, retrieval completeness, invalid inputs, and paired
metrics. A fictional fixture is included solely to test software behavior.

The planned retrieval controls are BM25, dense BGE-M3, reciprocal-rank-fusion
hybrid retrieval, contextual hybrid retrieval with reranking, and a gold
evidence oracle. Retrieval parameters (512/1,024-token chunks, 128-token
overlap, $k\in\{5,10,20\}$) are tuned on development data and then frozen.
The generation matrix will include models from at least three frontier
providers and one open-weight 8B--14B baseline, with temperature zero and
three repetitions. The full study has $60 \times 4 \times 5 \times 3 = 3{,}600$
generations before optional ablations.

The scientific contribution is not that one retriever dominates another.
Retrievers are controlled factors needed to separate an evidence-retrieval
failure from a failure to use evidence supplied to the model. The
\texttt{gold\_only} condition is the key oracle for this decomposition.

\section{Preliminary reproducibility audit}

We audited public artifacts from IPO Finance Agent at commit
\texttt{091f069}; the full revision is recorded in the released audit JSON.
This audit is descriptive:
it neither independently validates the benchmark's rubrics nor estimates model
accuracy. It establishes what can and cannot be treated as a gold-evidence
source when constructing \benchmark{}.

\begin{table}[t]
\caption{Public IPO Finance Agent question taxonomy at the audited commit. The
70-item public set is a useful seed pool, but its single-issuer concentration
motivates multi-filing construction for \benchmark{}.}
\label{tab:taxonomy}
\centering
\begin{tabular}{lr}
\toprule
Original category & Questions \\
\midrule
Quantitative & 28 \\
Disclosure & 21 \\
Forensic & 8 \\
Comparative & 6 \\
Governance & 4 \\
Modeling & 3 \\
\midrule
Total & 70 \\
\bottomrule
\end{tabular}
\end{table}

The public repository includes two automatically checked rubric artifacts
(Table~\ref{tab:rubrics}). The later pass reports higher mean overall quality,
but also contains 23 \texttt{stop} and 43 \texttt{repair} recommendations. No
record exposes a top-level evidence-identifier field. These upstream fields
must not be interpreted as human-validation error rates; they motivate a
separate, stable, claim-to-passage annotation layer.

\begin{table}[t]
\caption{Descriptive audit of upstream automatic rubric artifacts. Recommendation
labels and quality scores are reported exactly as stored upstream; they are not
independent human labels.}
\label{tab:rubrics}
\centering
\small
\begin{tabular}{lrrr}
\toprule
Artifact & Records & Mean quality & Recommendations \\
\midrule
Checked & 70 & 0.763 & 44 repair / 21 enrich / 5 stop \\
Seventh pass & 70 & 0.896 & 43 repair / 4 enrich / 23 stop \\
\bottomrule
\end{tabular}
\end{table}

One public Qwen result artifact contains 70 successful outputs. Sixty-eight
contain at least one URL (median one URL; maximum ten), but none contains a
Markdown inline citation or a stable evidence/chunk/passage ID token. This is a
format-level observation about one artifact, not a claim about factual
correctness, all models, or IPO Finance Agent as a whole. It does show why a
document-level URL cannot by itself establish citation completeness for an
atomic claim.

Finally, eight unit tests pass and the fictional fixture validates without schema
errors. Its smoke test distinguishes an intentionally evidence-sensitive policy
from an intentionally relevance-only policy. Those values are software
verification only; they are not included as benchmark results and must not be
used to support an empirical claim.

\section{Discussion and limitations}

The core risk is the validity of the human evidence sets. Minimality is
judgment-dependent, SEC filings can contain redundant disclosures, and
parametric model knowledge may permit a correct answer even after a controlled
removal. We mitigate rather than eliminate these risks through multiple MSES
paths, closed-book controls, frozen corpora, independent verification, and an
explicitly bounded claim: behavior is evaluated relative to the allowed corpus.

The financial domain is valuable because filings are long, consequential, and
traceable, but it limits generalization. The intervention protocol should later
be tested on legal, biomedical, and technical standards documents. Temporal
invalidity and contradictions are also important in finance; they are left as
post-core extensions rather than being folded into the first study.

This V0.1 is not ready for a NeurIPS submission. It has no completed
human-validated dataset and no real-model results. Its purpose is to make the
future experiment falsifiable, auditable, and difficult to overclaim. The
necessary next milestone is a double-annotated 20-item pilot meeting the
pre-specified quality gates.

\section{Conclusion}

\benchmark{} asks a narrower question than generic financial-agent quality:
does the agent change from supported answering to abstention when all known
necessary evidence paths are broken? The proposed paired design distinguishes
relevance from sufficiency, retrieval from evidence use, and citation presence
from claim-level support. The released V0.1 implementation and audit provide a
reproducible foundation; empirical conclusions await the planned human
annotation and model study.

\section*{Research transparency statement}

The implementation, fictional fixture, protocol, and reproducible upstream
audit are public. No private filing data, API keys, held-out questions, or
model results are included. The author list, venue, final dataset, and model
access details remain to be confirmed. A complete official checklist must be
generated for the final selected NeurIPS or workshop track.

{\small
\bibliographystyle{plainnat}
\bibliography{references}
}

\appendix
\section{Annotation record requirements}

Each record contains a question, a bounded allowed corpus, a reference answer,
atomic claims, evidence units, all known \mses{} paths, distractors,
answerability, temporal metadata, and adjudication metadata. The JSONL schema
and validation command are released with the code. Questions with unresolved
alternative evidence paths are excluded rather than forced into a
necessary-removal intervention.

\section{Reproducibility checklist for the future empirical study}

The final empirical release will provide filing hashes, data and prompt
versions, retrieval configuration, exact model IDs and access dates, run IDs,
costs, evaluator calibration, scripts for bootstrap intervals and tests, and
an error-analysis sample. This appendix is a planning commitment, not evidence
that the future experiment has already been performed.

\end{document}
